% ══════════════════════════════════════════════════════════════════
%  FICHE D'EXERCICES — Chapitre 3 : Les alcènes et les alcynes
% ══════════════════════════════════════════════════════════════════
\section{Exercices type Baccalauréat}
\textit{Données : $M(\ce{H})=1$, $M(\ce{C})=12$, $M(\ce{Br})=80$ \si{\gram\per\mol} ;
$V_m=\SI{24}{\litre\per\mol}$.}

\rubrique{Nomenclature, isomérie Z/E}

\exo{}\diff{1} Nommer : a)~$\ce{CH3-CH=CH-CH3}$ ; b)~$\ce{CH2=C(CH3)-CH3}$ ;
c)~$\ce{CH3-CH2-C#C-CH3}$ ; d)~$\ce{CH#C-CH2-CH3}$.

\exo{}\diff{1} Écrire les formules semi-développées de : a)~3-méthylbut-1-ène ;
b)~2-méthylbut-2-ène ; c)~pent-2-yne ; d)~4-méthylpent-2-yne.

\exo{}\diff{2} Écrire les formules semi-développées de tous les alcènes de formule
brute $\ce{C4H8}$, les nommer, et indiquer ceux qui présentent une isomérie $Z/E$
(représenter alors les deux stéréo-isomères).

\rubrique{Additions, tests d'insaturation}

\exo{}\diff{1} Écrire l'équation de l'addition : a)~de $\ce{H2}$ sur le but-1-ène ;
b)~de $\ce{Br2}$ sur l'éthène ; c)~de $\ce{HCl}$ sur le propène (énoncer la règle de
Markovnikov et donner le produit majoritaire).

\exo{}\diff{2} On fait réagir du propène avec de l'eau (en milieu acide).
\begin{enumerate}[label=\arabic*)]
\item Écrire l'équation de l'hydratation.
\item Donner les deux produits possibles, les nommer et préciser le produit
majoritaire (règle de Markovnikov).
\end{enumerate}

\exo{}\diff{2} Un alcène $A$ décolore l'eau de brome. L'addition de dibrome sur
\SI{2.8}{\gram} de $A$ consomme exactement \SI{0.05}{\mol} de $\ce{Br2}$.
\begin{enumerate}[label=\arabic*)]
\item Déterminer la masse molaire de $A$ et sa formule brute.
\item Donner les formules semi-développées possibles de $A$ et leurs noms.
\end{enumerate}

\exo{}\diff{2} Comment distinguer expérimentalement un alcane d'un alcène ? Décrire
le test à l'eau de brome et écrire l'équation correspondante pour le but-2-ène.

\rubrique{Polymérisation, alcynes}

\exo{}\diff{2} L'éthylène (éthène) se polymérise en polyéthylène.
\begin{enumerate}[label=\arabic*)]
\item Écrire l'équation de polymérisation (motif et indice $n$).
\item Calculer le degré de polymérisation $n$ d'un échantillon de masse molaire
moyenne $M = \SI{84000}{\gram\per\mol}$.
\end{enumerate}

\exo{}\diff{2} L'éthyne (acétylène) $\ce{C2H2}$ donne, par addition successive de
dihydrogène, l'éthène puis l'éthane. Écrire les deux équations et nommer les produits.

\exo{}\diff{3} L'hydratation de l'éthyne (catalyse) conduit à un composé $B$ qui
donne un test positif à la 2,4-DNPH et au réactif de Tollens. Identifier $B$, écrire
l'équation de sa formation, et justifier les résultats des tests.

\rubrique{Problèmes de synthèse — type Baccalauréat}

\sujetbac[d'après Baccalauréat C, Cameroun]
Un hydrocarbure $A$ a pour densité de vapeur par rapport à l'air $d = 1{,}93$. Il
décolore l'eau de brome et fixe le dibrome dans le rapport molaire $1{:}1$.
\begin{enumerate}[label=\arabic*)]
\item Déterminer la masse molaire de $A$ et montrer qu'il s'agit d'un alcène de
formule $\ce{C4H8}$.
\item L'hydratation de $A$ ne donne qu'un seul alcool, qui est un alcool tertiaire.
Identifier $A$, le nommer, et écrire l'équation de son hydratation.
\item Écrire l'équation de l'addition de $\ce{HBr}$ sur $A$ et nommer le produit majoritaire.
\end{enumerate}

\sujetbac[Identification par hydrogénation]
L'hydrogénation complète de \SI{2.0}{\gram} d'un hydrocarbure insaturé $X$ consomme
\SI{1.8}{\litre} de dihydrogène ($V_m=\SI{24}{\litre\per\mol}$) et fournit un alcane $Y$.
\begin{enumerate}[label=\arabic*)]
\item Déterminer le nombre de moles de $\ce{H2}$ fixées par mole de $X$ (on donne
$M(X)=\SI{54}{\gram\per\mol}$).
\item En déduire le nombre d'insaturations de $X$ et proposer une formule brute.
\item Proposer une formule semi-développée pour $X$ (un diène ou un alcyne) et écrire
l'équation de son hydrogénation totale.
\end{enumerate}
